\documentclass[11pt]{article}
\usepackage[margin=1in]{geometry}
\usepackage{amsmath,amssymb,mathtools}
\usepackage{hyperref}
\usepackage{pdfpages}
\setlength{\parindent}{0pt}
\setlength{\parskip}{0.55em}
\newcommand{\R}{\mathbb R}
\begin{document}
\begin{center}
{\Large\bfseries Literature resolution of the sum-bifunction question}\\[3pt]
{\large arXiv:1105.5550, Remark 7}
\end{center}

\textbf{Status.} Literature already answered by a later general
necessary-and-sufficient characterization.  The application below is an
inference; the later paper does not explicitly label its theorem as an answer
to Remark 7.

\textbf{The source question.}  Let $X$ be reflexive, let $C\subset X$ be
closed and convex, and let $F,G:C\times C\to\R$ be monotone, vanish on the
diagonal, and have convex lower-semicontinuous second-variable slices.  Under
$0\in\operatorname{sqri}(C-C)$, Proposition 18 proves
\[
 A^F+A^G=A^{F+G}.
\]
Remark 7 asks what additional hypothesis, no stronger than the preceding
standard sufficient hypotheses, guarantees maximality of $F+G$.

\textbf{Later answer.}  Hadjisavvas, Jacinto, and Mart\'inez-Legaz prove in
Theorem 2 of \emph{Some Conditions for Maximal Monotonicity of Bifunctions}
(published online 2015, DOI
\href{https://doi.org/10.1007/s11228-015-0343-6}{10.1007/s11228-015-0343-6})
that, for every monotone bifunction $Q$ with convex lower-semicontinuous
second-variable slices, maximal monotonicity is equivalent to three other
conditions.  One is pointwise maximality of the canonical support
bifunction $G_{A^Q}$.  Another is the exact comparison condition
\begin{equation}
\begin{gathered}
 \text{for every admissible PMMB }K\text{ with finite Fitzpatrick transform,}\\
 \exists x_K:\qquad Q(x_K,y)+K(x_K,y)\ge0\qquad(\forall y).
\end{gathered}
\tag{1}
\end{equation}
Their fourth clause is an equivalent perturbed-equilibrium solvability
condition for every $(x_0,x_0^*)\in X\times X^*$.

To apply the theorem, put $Q=F+G$ and extend it from $C\times C$ by
\[
 \widehat Q(x,y)=
 \begin{cases}
 Q(x,y),&x,y\in C,\\
 +\infty,&x\in C,\ y\notin C,\\
 -\infty,&x\notin C.
 \end{cases}
 \tag{2}
\]
Then $\widehat Q$ is monotone, its second-variable slices are convex and
lower semicontinuous, and $A^{\widehat Q}=A^Q$.  Thus Theorem 2 applies
directly.  Conditions (1) and its equivalent clauses are exact additions to
Proposition 18: they are necessary and sufficient, and hence cannot be
stronger than any valid sufficient condition from the preceding remark.

\textbf{Why something must be added.}  The original assumptions alone do
not imply maximality.  Take $X=C=\R$, $G=0$, and
\[
 a(x)=\begin{cases}0,&x\le0,\\1,&x>0,\end{cases}
 \qquad F(x,y)=a(x)(y-x).
\]
The function $a$ is nondecreasing, so $F$ is monotone; every $F(x,\cdot)$ is
affine and continuous, and $0\in\operatorname{sqri}(C-C)$.  Nevertheless
$A^F(x)=\{a(x)\}$ is not maximal: each $(0,t)$, $0<t\le1$, is monotonically
related to its graph but absent from it.

\textbf{Scope.}  The 2015 theorem resolves the broad request by an exact
characterization.  It does not turn the answer into a single elementary
semicontinuity assumption, so searches for especially simple sufficient
conditions remain meaningful but are no longer the same open question.

\newpage
\thispagestyle{empty}
\includepdf[pages=10,pagecommand={\thispagestyle{empty}},fitpaper=true]{source_paper.pdf}
\includepdf[pages=7,pagecommand={\thispagestyle{empty}},fitpaper=true]{supporting_paper_2015.pdf}
\end{document}
